% Replacement for the rational-source example and its obstruction proof.
% Inputs: parent, change, normal, split, lmis, three, and prop:degree
% from canonical_completion.tex.  No additional scientific input is used.
Consider the C1 source $\lambda=n=0$ with
\begin{equation}
 (d,m,C,A,D,\alpha,\beta)
 =\left(359,\frac{75}{16},60,416,
              \frac{41}{2},\frac{831}{2},\frac{151}{2}\right).
 \label{simple:source}
\end{equation}
All six event weights are strictly positive. Dividing this tuple by $359$
gives the same inequality with $d=1$; the scaling in
\eqref{simple:source} keeps the following calculations short.

\paragraph{Validity of the source.}
Let $M_*(u,v)$ be the symmetric matrix of the full parent
\eqref{parent} in $(W,X,Y,Z)$. Put $\xi=4u^2$ and $\eta=v^2/4$.
Direct expansion gives, for $uv\ne0$,
\begin{align}
 \det M_*(u,v)
 &=28800D_u^2D_v^2u^2v^2\,\mathcal L(\xi,\eta),
 \label{simple:determinant}\\
 \mathcal L(\xi,\eta)
 &=3757645(\xi+\xi^{-1}-2)
   +14022297(\eta+\eta^{-1}-2)\notag\\
 &\quad+3732435(\eta/\xi+\xi+\eta^{-1}-3)
   +8179670(\xi\eta+\xi^{-1}+\eta^{-1}-3).
 \label{simple:balanced}
\end{align}
Each bracket is nonnegative by the two- or three-term arithmetic--geometric
mean inequality: its positive monomials have product one.
Thus the full state determinant is nonnegative.

The $Z=0$ principal block $B_*$ is strictly positive for $u\ne0$.
Indeed, its three positive $m,C,A$ terms give
\[
 B_*\succeq
 T^T\operatorname{diag}(4m,4CD_v,4AD_uD_v)T,
 \qquad
 T=\begin{pmatrix}v&-1&uv\\u&0&-1\\0&0&1\end{pmatrix},
 \qquad \det T=u.
\]
All remaining terms restricted to $Z=0$ are nonnegative.
The scalar Schur complement therefore proves $M_*\succeq0$ whenever
$uv\ne0$. Polynomial continuity gives both finite axes. Homogenizing
the parent to bidegree $(2,2)$ in the two projective angle pairs gives
their infinite endpoints as well. Both angle signs are included, and
separating real and imaginary state amplitudes gives the same conclusion
for complex states. This proves complete validity of the source.

An exact contact is
\begin{equation}
 u_*=\frac12,\qquad v_*=2,\qquad
 (W_*,X_*,Y_*,Z_*)=\left(2,-\frac14,\frac14,1\right).
 \label{simple:physical-contact}
\end{equation}
The full state equation $M_*(u_*,v_*)\psi_*=0$ and both angle
derivatives of the parent vanish exactly. In the change of coordinates
\eqref{change}, the homogeneous state $10\psi_*$ is
\begin{equation}
 U_*=\frac14,\qquad
 a_* =\binom{17}{2},\qquad b_* =\binom{12}{-18}.
 \label{contact}
\end{equation}

\begin{theorem}\label{thm:obstruction}
The valid source \eqref{simple:source} has no polynomial selector satisfying
\eqref{lmis} for every real $v$. Equivalently, its polynomial
$\Phi_c(u^2,v;a,b)$ is not a sum of polynomial squares in $(u,v,a,b)$.
\end{theorem}
\begin{proof}
By Proposition~\ref{prop:degree}, a polynomial selector can be chosen as
\eqref{three}. At the contact, both nonnegative summands in
\eqref{split} vanish. Since a positive semidefinite matrix kills every
vector on which its quadratic form vanishes, this forces
\begin{equation}
 J(2)b_*=-P(2)a_*/U_*=\binom{-4800}{-38400},\qquad
 J(2)^Ta_*=-U_*N(2)b_*=\binom{1200}{9600}.
 \label{kernel}
\end{equation}
For completeness, the two matrices here are
\[
 P(2)=\begin{pmatrix}75&-75/2\\-75/2&20475/4\end{pmatrix},
 \qquad
 N(2)=\begin{pmatrix}2000&1600\\1600&3200\end{pmatrix}.
\]
Writing $J(2)=\bigl(\begin{smallmatrix}4p&0\\2r&4s\end{smallmatrix}\bigr)$,
the first coordinate of the first equation, and the two coordinates of
the second, give
\[
 48p=-4800,\qquad 68p+4r=1200,\qquad 8s=9600.
\]
The remaining equation $24r-72s=-38400$ is then satisfied.
Consequently the unique possible polynomial selector is
\begin{equation}
 p=-100,\qquad r=2000,\qquad s=1200.
 \label{forced}
\end{equation}
At $v=1$, the principal submatrix of $\mathcal G_J$ on coordinates
$(a_2,b_1,b_2)$ is therefore
\begin{equation}
 \begin{pmatrix}
 8235/4&2000&1200\\
 2000&1280&800\\
 1200&800&800
 \end{pmatrix},
 \qquad \det=-412640000<0.
 \label{minor}
\end{equation}
This contradicts $\mathcal G_J(1)\succeq0$.
Positive scaling of the source scales $P,N,E,J$ by the same factor;
thus the normalized $d=1$ source has selector parameters
$(-100,2000,1200)/359$ and the negative minor
$-412640000/359^3$.

Finally, a polynomial sum of squares for $\Phi_c(u^2,v;a,b)$ can use
roots of degree at most two in $u$, at most one in $v$, and homogeneous
linear degree in the state coordinates: higher leading degrees cannot
cancel in a sum of real squares. Averaging under $u\mapsto-u$ and using
the constant and leading state supports, as in
Theorem~\ref{thm:selector}, then produces a polynomial selector.
Conversely, a polynomial selector gives polynomial matrix factorizations
of its two positive univariate matrices and hence a polynomial sum of
squares through \eqref{split}. This proves the stated equivalence.
\end{proof}

This obstruction is specific to polynomial selectors in this completion.
It excludes neither an algebraic selector nor a certificate after an
additional fixed testing multiplier, and does not exclude a finite
semidefinite lift.
